\documentclass[11pt,a4paper]{article}
\usepackage[top=1.7cm,bottom=1.5cm,left=1.5cm,right=1.5cm,headsep=8pt,footskip=16pt]{geometry}
\usepackage{enumitem}
\usepackage{amssymb}
\usepackage{array}
\usepackage[table]{xcolor}
\usepackage{fancyhdr}
\usepackage{lastpage}
\usepackage{multicol}
\usepackage{polyglossia}
\setdefaultlanguage[calendar=gregorian, numerals=maghrib]{arabic}
\setotherlanguage{english}
\IfFontExistsTF{Amiri}{%
  \newfontfamily\arabicfont[Script=Arabic, Scale=1.08]{Amiri}
}{%
  \newfontfamily\arabicfont[Script=Arabic, Scale=1.08, Path=../fonts/, Extension=.ttf,
    UprightFont=*-Regular, BoldFont=*-Bold, ItalicFont=*-Italic,
    BoldItalicFont=*-BoldItalic]{Amiri}
}
\IfFontExistsTF{AlBattar}{%
  \newfontfamily\khatt[Script=Arabic, Scale=1.15]{AlBattar}
}{%
  \newfontfamily\khatt[Script=Arabic, Scale=1.15, Path=../fonts/,
    BoldFont=ae_AlBattar.ttf]{ae_AlBattar.ttf}
}
\renewcommand{\headrulewidth}{0.8pt}
\renewcommand{\footrulewidth}{0.8pt}
\fancyhf{}
\rhead{\khatt المدرسة العليا للأساتذة ببوزريعة}
\lhead{\khatt ملحقة بومرداس}
\lfoot{\khatt الأستاذ: د. بن عبد الرزاق يوسف}
\cfoot{\thepage/\pageref{LastPage}}
\pagestyle{fancy}
\parindent 0pt
\setlength{\columnsep}{0.6cm}
\setlength{\columnseprule}{0.4pt}
\setlist[enumerate,1]{label=\textbf{\arabic*.}, leftmargin=*, itemsep=5pt, topsep=2pt}
\begin{document}

\linespread{1.08}\selectfont
\begin{center}
{\LARGE\khatt\textbf{الاختبار الخامس عشر في الذكاء الاصطناعي}}\\[4pt]
{\normalsize وضعيات إشكالية ــ أسئلة متعددة الاختيارات ــ الفصول من 1 إلى 6}
\end{center}
\vspace{-2pt}\rule{\textwidth}{0.6pt}\\[2pt]
\textbf{التعليمة:} يتضمن الاختبار خمس وضعيات إشكالية. اقرأ كل وضعية بعناية، ثم اختر التحليل أو التصرف الأنسب
من بين الاختيارات الخمسة؛ لكل وضعية إجابة صحيحة واحدة فقط.
اكتب \textbf{حرف} الإجابة الصحيحة في جدول الإجابات في آخر الاختبار.
لكل وضعية 4 نقاط (المجموع 20 نقطة).
\vspace{4pt}

\noindent
{\large\khatt\textbf{الوضعية 1}}\hfill{\small(الفصل 1)}\par\nopagebreak\vspace{3pt}
\setlength{\fboxsep}{7pt}\colorbox{gray!10}{\parbox{\dimexpr\linewidth-14pt}{تفكر أستاذة في تحويل الأداة التوليدية إلى «مساعد» دائم لتلاميذ السنة الرابعة متوسط يسألونه في البيت عن دروس العربية. ويقترح زميل أن يحل المساعد محل مراجعة الأستاذة للدروس في القسم. والتلاميذ متفاوتون في المستوى وفي الوصول إلى الإنترنت. وستعرض الفكرة على مجلس الأساتذة الأسبوع القادم. ما التقدير الأنسب؟}}\par\nopagebreak
\begin{itemize}[leftmargin=1.0cm, labelsep=6pt, itemsep=3pt, topsep=4pt, parsep=0pt]
  \item[\textbf{أ)}] يُعتمد المساعد للتلاميذ المتفوقين وحدهم، لأنهم يفهمون إجابات الأداة، أما الضعاف فيُمنعون منه، ويُلغى التدريب على التحقق لأن المتفوقين لا يحتاجون إليه.
  \item[\textbf{ب)}] يمكن أن يكون مساعدًا مكمّلًا بعد تدريب التلاميذ على صياغة الأسئلة والتحقق من الإجابات، مع إبقاء المراجعة في القسم وتوفير بدائل لمن لا يملك الوصول، لأن الأداة قد تخطئ ولا تعوض الأستاذ.
  \item[\textbf{ج)}] يُعتمد المساعد مكمّلًا مع إبقاء المراجعة في القسم، دون تدريب التلاميذ على التحقق، لأن الأداة المخصصة للتعليم لا تخطئ في دروس المرحلة المتوسطة أبدًا.
  \item[\textbf{د)}] يحل المساعد محل المراجعة في القسم تمامًا، لأن الأداة تجيب عن كل الأسئلة دون أخطاء، والتلاميذ الذين لا يملكون الإنترنت يمكنهم الاعتماد على زملائهم في نقل الإجابات، فيتفرغ الأستاذ لمهام أخرى أكثر أهمية.
  \item[\textbf{هـ)}] يُرفض المساعد نهائيًا، لأن أي استعمال للأدوات التوليدية خارج القسم يضر بالتعلم حتمًا، ولا يمكن تدريب التلاميذ على التحقق من الإجابات في هذه المرحلة العمرية، ولو وُفِّرت لهم كل الوسائل والتوجيهات.
\end{itemize}
\par\vspace{6pt}
\noindent
{\large\khatt\textbf{الوضعية 2}}\hfill{\small(الفصل 2)}\par\nopagebreak\vspace{3pt}
\setlength{\fboxsep}{7pt}\colorbox{gray!10}{\parbox{\dimexpr\linewidth-14pt}{جرّبت أستاذة أداتين: الأولى نموذج ضخم عام، والثانية نموذج أصغر دُرِّب خصيصًا على نصوص نحوية عربية محققة. في أسئلة الإعراب الدقيقة تفوقت الثانية، وفي كتابة النصوص الإبداعية تفوقت الأولى. وتريد أن تفهم سبب هذا التفاوت. وتريد أن تختار الأداة المناسبة لكل مهمة في عملها اليومي. ما التفسير الأدق؟}}\par\nopagebreak
\begin{itemize}[leftmargin=1.0cm, labelsep=6pt, itemsep=3pt, topsep=4pt, parsep=0pt]
  \item[\textbf{أ)}] التفاوت سببه أن الأستاذة كتبت أسئلة الإعراب بأسلوب أفضل للنموذج الثاني، فالبيانات والحجم لا يؤثران في الأداء، والطلب وحده يحدد جودة الإجابة في كل نموذج، وهذا ما تثبته تجارب كثيرة في المدارس.
  \item[\textbf{ب)}] النموذج الأصغر أفضل في كل شيء بالضرورة، لأن النماذج الصغيرة أدق دائمًا، وتفوق الأكبر في الإبداع سببه أن الإبداع لا يحتاج إلى أي تعلّم من البيانات.
  \item[\textbf{ج)}] أداء النموذج يتأثر بحجمه وبطبيعة بياناته معًا؛ النموذج المتخصص تعلّم أنماط الإعراب من بيانات مركّزة عالية الجودة، والعام تعلّم من تنوع واسع يخدم الإبداع، فلا يتفوق نموذج في كل المهام.
  \item[\textbf{د)}] النموذج المتخصص تعلّم الإعراب من بيانات مركّزة، وهذا يعني أنه لا يخطئ في الإعراب أبدًا، فتعتمد إجاباته الإعرابية دون مراجعة، وتراجع النموذج العام وحده، فتوفر بذلك وقتها في مراجعة الإجابات الإعرابية.
  \item[\textbf{هـ)}] النموذج الأكبر أفضل في كل شيء بالضرورة، وتفوق الأصغر في الإعراب مصادفة لا تتكرر، لأن الحجم وحده يحدد جودة الأداء في كل المهام اللغوية.
\end{itemize}
\par\vspace{6pt}
\noindent
{\large\khatt\textbf{الوضعية 3}}\hfill{\small(الفصل 3)}\par\nopagebreak\vspace{3pt}
\setlength{\fboxsep}{7pt}\colorbox{gray!10}{\parbox{\dimexpr\linewidth-14pt}{يعدّ أستاذ أسبوعيًا تمارين لأقسام مختلفة، ويعيد كتابة الطلب كل مرة من البداية، فتتفاوت جودة النتائج من أسبوع إلى آخر. واقترح زميل أن يبني «نموذج طلب» ثابتًا يعيد استعماله. وقد لاحظ أن أفضل نتائجه جاءت عندما قدّم أمثلة واضحة. ما المكونات الأنسب لهذا النموذج حتى يحقق الاتساق؟}}\par\nopagebreak
\begin{itemize}[leftmargin=1.0cm, labelsep=6pt, itemsep=3pt, topsep=4pt, parsep=0pt]
  \item[\textbf{أ)}] قالب يتضمن مثالًا نموذجيًا فقط دون أي معلومات عن المستوى أو الدرس، لأن المثال وحده كافٍ لكل الأقسام، والأداة تستنتج المستوى والدرس منه في كل مرة.
  \item[\textbf{ب)}] جملة ثابتة قصيرة مثل «اكتب تمارين جيدة للدرس»، لأن القالب القصير أسهل في إعادة الاستعمال، والأداة تتذكر تفضيلاته من الأسابيع السابقة فتحقق الاتساق تلقائيًا.
  \item[\textbf{ج)}] قالب طويل يتضمن كل الدروس والمستويات معًا في طلب واحد، لأن جمع كل شيء يضمن الاتساق، ولا حاجة إلى خانات متغيرة أو مثال نموذجي أو مراجعة النتائج.
  \item[\textbf{د)}] قالب ثابت يتضمن الدور والمستوى والدرس وشكل الإخراج، لكنه لا يغيّر أي عنصر منه أبدًا مهما تكررت المشكلات، لأن تعديل القالب يُفقده قيمته في تحقيق الاتساق.
  \item[\textbf{هـ)}] قالب ثابت يتضمن: الدور، والمستوى، والدرس المستهدف، والمهارة المطلوبة، وعدد الأسئلة، وشكل الإخراج، ومثالًا نموذجيًا، مع خانات يغيّرها كل أسبوع، ثم يراجع النتائج ويحسّن القالب تدريجيًا.
\end{itemize}
\par\vspace{6pt}
\noindent
{\large\khatt\textbf{الوضعية 4}}\hfill{\small(الفصل 5)}\par\nopagebreak\vspace{3pt}
\setlength{\fboxsep}{7pt}\colorbox{gray!10}{\parbox{\dimexpr\linewidth-14pt}{يريد طالب ترجمة فقرة من كتاب فرنسي في النقد الأدبي للاستشهاد بها في مذكرته، فترجمها بأداة ذكاء اصطناعي، وجاءت الترجمة سلسة، لكنه لاحظ أن مصطلحين نقديين تُرجما بألفاظ غير مستعملة في المراجع العربية المتخصصة. وموعد تسليم المذكرة بعد شهر، ومشرفه يدقق في المصطلحات النقدية تدقيقًا شديدًا. كيف يتصرف تصرفًا منهجيًا؟}}\par\nopagebreak
\begin{itemize}[leftmargin=1.0cm, labelsep=6pt, itemsep=3pt, topsep=4pt, parsep=0pt]
  \item[\textbf{أ)}] يقابل الترجمة بالأصل، ويرجع إلى معاجم المصطلحات النقدية والمراجع العربية لاختيار المقابل المعتمد، ويذكر المصطلح الفرنسي بين قوسين، ويوثق الأصل مع الإشارة إلى أن الترجمة من إعداده.
  \item[\textbf{ب)}] يطلب ترجمة ثانية بأداة أخرى، فإن اتفقت الأداتان على المصطلحين اعتمدهما، لأن اتفاق أداتين دليل على أن المقابل هو المعتمد في المراجع العربية المتخصصة، ثم يوثّق الترجمة في الهامش باسم الأداتين.
  \item[\textbf{ج)}] يعتمد الترجمة كما هي، لأن الأداة اختارت أحدث المقابلات، والمراجع العربية متأخرة في المصطلح، ويوثّق الترجمة باسم الأداة بدل الكتاب الفرنسي الأصلي.
  \item[\textbf{د)}] يستبدل المصطلحين بأي مرادف يراه مناسبًا دون الرجوع إلى المراجع، ويحذف الإحالة إلى الأصل الفرنسي، لأن الترجمة تجعل الفقرة نصًا جديدًا لا يحتاج إلى توثيق.
  \item[\textbf{هـ)}] يتخلى عن الاستشهاد بالفقرة، لأن أي ترجمة آلية لنص نقدي غير مقبولة علميًا، حتى لو روجعت مصطلحاتها وقوبلت بالأصل ووُثِّقت توثيقًا كاملًا.
\end{itemize}
\par\vspace{6pt}
\noindent
{\large\khatt\textbf{الوضعية 5}}\hfill{\small(الفصل 6)}\par\nopagebreak\vspace{3pt}
\setlength{\fboxsep}{7pt}\colorbox{gray!10}{\parbox{\dimexpr\linewidth-14pt}{طلبت أستاذة من أداة توليد صور إنتاج إنفوغرافيك عن «مراحل كتابة المقال» بخمس مراحل مرقمة. جاء التصميم جذابًا، لكن الترتيب البصري للمراحل غير واضح، وبعض الأرقام مكررة، والنصوص العربية ناقصة الحروف. وتريد توزيعه مطبوعًا على التلاميذ. وعدد التلاميذ يتجاوز مئة وعشرين تلميذًا. ما الخطة الأنسب؟}}\par\nopagebreak
\begin{itemize}[leftmargin=1.0cm, labelsep=6pt, itemsep=3pt, topsep=4pt, parsep=0pt]
  \item[\textbf{أ)}] تطلب من الأداة كتابة النصوص بالإنجليزية داخل الإنفوغرافيك، لأن الأداة تتقنها، ثم توزعه على التلاميذ ليترجموه، فيصبح النشاط تمرينًا لغويًا إضافيًا.
  \item[\textbf{ب)}] تطبعه كما هو، لأن الجاذبية البصرية أهم من دقة الترتيب، والتلاميذ يستطيعون استنتاج الترتيب الصحيح للمراحل وإكمال الحروف الناقصة بأنفسهم.
  \item[\textbf{ج)}] تعيد الطلب حتى يظهر إنفوغرافيك بلا أخطاء ظاهرة، ثم تطبعه دون تدقيق الأرقام والنصوص، لأن غياب الأخطاء الظاهرة يعني أن المحتوى صحيح بالضرورة، ثم توزعه على التلاميذ في أقرب حصة.
  \item[\textbf{د)}] تحدد المراحل ونصوصها بنفسها، وتستعمل الأداة لاقتراح التصميم أو الرموز فقط، ثم تركّب الإنفوغرافيك في برنامج تصميم تضبط فيه الترتيب والأرقام والنصوص، وتراجعه قبل الطباعة.
  \item[\textbf{هـ)}] تتخلى عن الإنفوغرافيك نهائيًا، لأن أدوات توليد الصور لا يمكن أن تساعد في أي جزء من التصميم التعليمي، حتى في اقتراح الرموز والألوان والتخطيط العام.
\end{itemize}
\par\vspace{6pt}
\begin{center}
{\khatt\textbf{جدول الإجابات}}\\[3pt]
\renewcommand{\arraystretch}{1.6}
\begin{tabular}{|>{\centering\arraybackslash}m{1.9cm}|*{5}{>{\centering\arraybackslash}m{1.2cm}|}}
\hline
\textbf{رقم السؤال} & \textbf{1} & \textbf{2} & \textbf{3} & \textbf{4} & \textbf{5} \\ \hline
\textbf{الإجابة} &  &  &  &  &  \\[8pt] \hline
\end{tabular}

\vspace{8pt}
\renewcommand{\arraystretch}{2.2}
\begin{tabular}{|>{\centering\arraybackslash}m{7.5cm}|>{\centering\arraybackslash}m{3.2cm}|>{\centering\arraybackslash}m{3.2cm}|}
\hline
\textbf{الاسم واللقب} & \textbf{الفوج} & \textbf{العلامة} \\ \hline
 & & \hfill /20 \\ \hline
\end{tabular}
\end{center}
\end{document}
